% !TeX root = ../papers/nc-p3.tex
\section{特異・可約な種数一の四次部分曲線と印の移動}
\label{sec:genus-one-point-marks}
各辺の左右の射影ファイバーが空か一点、像 $X$ が一次元である有限Frobenius窓を扱う。
$X$ の純四次部分曲線 $Y$ がGorensteinで、$H^0(\mathcal O_Y)=k$、
$\omega_Y=\mathcal O_Y$ を満たすときこれを印として使う。重複成分も許す。

\begin{lemma}\label{lem:genus-one-quartic-mark}
この印は各位置で二つの二次曲面の完全交叉となる。逆に一位置で型 $(2,2)$ の完全交叉部分曲線なら
全ての位置で上の印となる。相異なる印が同じ既約成分を共有することを排除しない。
印上の $L_i$ の積の全切断環 $D(Y)$ は一次生成で、
長さ $n>0$ の各成分の次元は $4n$ である。
\end{lemma}
\begin{proof}
各成分の次数は定理\ref{thm:point-graph-curve-degree}で一定であり、局所長を掛けて足せば
全位置で総次数4となる。真の純部分曲線 $Z\subsetneq Y$ の双対化層は
$\operatorname{Hom}_Y(\mathcal O_Z,\mathcal O_Y)$ なので大域切断は零である。
従って $h^1(\mathcal O_Z)=0$。多重次数零の可逆層 $M$ については
$H^0(M)\ne0$ なら $M=\mathcal O_Y$ であることも示す。
非零切断の像はある純部分曲線 $Z$ の構造層である（可逆層内に有限台の部分はない）。
$Z$ が真なら $M|_Z$ は自明である。実際、曲線の各被約成分で次数を取る写像の核は連結であり、
$h^1(\mathcal O_Z)=0$ ならこの核は零である。具体的には被約曲線の正規化の単元の完全列で、
各特異点の貼合せの群はトーラスと単冪群の拡大、滑らかな各成分の次数零Picard群は連結である。
冪零部分を戻すときも $1+\mathcal N^j$ の逐次商は加法群である。
これらが次数零部分の連結性を与え、その接空間は $H^1(\mathcal O_Z)$ である。
すると $\operatorname{Hom}_Y(\mathcal O_Z,M)=H^0(\omega_Z\otimes M|_Z)=0$ に反する。
$Z=Y$ なら切断の有限余核はEuler標数差が零なので消える。

一般の超平面切断は非零因子で非空な零因子を持つ。
従って正の長さの偏極積の逆束は $\mathcal O_Y$ に入り大域切断が零である。
双対性とRiemann--Rochにより長さ $n>0$ の積は $H^1=0,h^0=4n$ となる。
平面内の純四次曲線は算術種数3なので各埋込みは非退化で、完全線型系になる。
二次式は少なくとも二本ある。独立な二本が $\ell a,\ell b$ と共通平面因子を持つとする。
$\mathcal O_Y$ のイデアル $(\ell)$ は直線 $a=b=0$ 上の $\mathcal O(-1)$ の商であり、
非零で、その台の純性から同じ $\mathcal O(-1)$ に等しい。
従って $Y\cap\{\ell=0\}$ は次数3、Euler標数0である。平面内の純三次部分のEuler標数は0
なので孤立・埋込み点はなく、真の純部分曲線に $h^1=1$ を持つことになり矛盾する。
ゆえに二式は完全交叉であり、次数一致と純性で $Y$ と等しい。逆向きは随伴と完全交叉の連結性
（構造層のKoszul解消で $h^0=1$）である。

$L_i$ の一般の二切断は基点自由である。鉛筆の完全列を次の束 $L_{i+1}$ でひねると
切断乗法の余核は $H^1(L_{i+1}/L_i)$ の商である。多重次数零の先の議論と双対性により
$L_i\ne L_{i+1}$ なら零である。一致する場合は二次完全交叉の射影正規性を使う。
長さ3以上では同じ鉛筆列に正の尾部がさらに掛かる。
障害群の双対は多重次数零の束に正の偏極積の逆を掛けた束の切断である。
一般の正の切断を掛ければ多重次数零の束に単射し、後者の切断は零または定数なので障害は消える。
従って全切断乗法が全射となる。
\end{proof}

\begin{theorem}\label{thm:finite-orbit-quartic-marks}
四移動が印 $Y$ に持つ軌道 $\mathcal M$ が有限であるとする。
$W_i(Z)=\ker(A_{i,i+2}\to D(Z)_{i,i+2})$ は印 $Z$ ごとの二次元空間である。
次の条件が全次数の標準Koszul左右AS正則性と同値である。
軌道上に一つの置換 $\pi$ があり、各 $Z\in\mathcal M$ と各位置で
\[
 \ker\bigl(W_i(Z)\otimes W_{i+2}(\pi Z)\longrightarrow A_{i,i+4}\bigr)
\]
が階数2のテンソルを含む。位置は $4|\mathcal M|$ を法として調べればよい。
特に被約な四次完全交叉部分曲線を一つでも含む枝と、
重複を許す四次完全交叉が一次元部分の全体を占める枝にこの有限判定が適用される。
\end{theorem}
\begin{proof}
一つの固定した印で定理\ref{thm:elliptic-quartic-as-test}の束複体を評価する。
局所分裂、全ての正の偏極積に対する $H^1$ 消滅、$\omega_Z=\mathcal O_Z$ と
文字の完全線型系だけが使われている。従って同じ証明で左右とも第一Torは次数2に二つ、
第二Torは次数4に一つの単純の商である。標準の場合はこれで長さ2の左右解消を得る。
$4|\mathcal M|$ 頂点に折り返せば $D(Z)$ はこの代数の商である。
同定 $R^2\Gamma_A(D(Z))'\simeq\underline{\operatorname{Ext}}^2_A(D(Z),A)$ に付随するねじれは
標準四移動またはその逆（および次数スカラー）であり、全ての始点・印に共通である。
従って最後の右syzygyの二係数は今度は四移動した印の二次核を張る。
実際、ねじれたイデアルの左零化と二次元性から従う。
軌道でこのねじれを置換 $\pi$ として使えば必要性を得る。

逆に任意の $i,Z$ を固定する。偶奇二列の印 $Z_j$ を
$Z_i=Z,\ Z_{j+2}=\pi Z_j$ で選ぶ。他方の偶奇列の初期値は任意でよい。
必要なら周期の公倍数に展開する。矢 $W_j(Z_j)$ と仮定の非退化関係による抽象部分表示 $P$
は二変数多項式型である。始点 $j$ の矢で割ると $e_jD(Z_j)$、終点 $j+2$ の矢で割ると
$D(Z_j)e_{j+2}$ となる。どちらも同じHilbert関数を持つ。従って基底の持上げと
左からの $P$ の語で $A$ が張られ、逆向きにも張られる。
始点 $j$ の二次核の関係は四次に高々一元で生成されるので、仮定の関係がこの始点の全関係を生成する。
定理\ref{thm:elliptic-quartic-as-test}と同じ固定終点のEuler不等式より
$H_A=(1-z)^{-4}$、両側の張る集合が基底となる。特に $e_i D(Z)$ は $P$ 単純の誘導で
$\underline{\operatorname{Ext}}_A^q(e_iD(Z),A)=0$ ($q<2$) を得る。
同じ印 $Z$ を固定したまま各始点 $i$ でこの構成を行えることが重要である。
全始点の消滅を合わせれば、この一つの商 $A\to D(Z)$ の環変更で単純からの消滅を得る。
反対側では所定の終点の二つ手前の印を指定する。Frobenius双対とEuler式で正則性が従う。

被約な印は有限個の既約成分の和集合に限られる。
また一次元部分の全部を占める純な印は、有限台の部分を除いた商として一意である。
したがって最後の二場合はいずれも有限軌道である。
\end{proof}

\paragraph{印を指定しない場合。}
被約な印は射影像の被約既約成分の和をすべて調べ、純四次かつ二つの二次式の完全交叉になるものを残せば
有限に列挙できる。候補の同型は、二連続窓と位置 $0,1,2$ の偏極および選んだ印上の
曲線同型で比較できる。具体的には印上で三位置とも完全線型系なので、偏極を保つ曲線同型の引戻しが
二窓の関係を移せば、系\ref{cor:two-window-extension}で十分である。一つ目で固定した印を
二つ目の有限リストのいずれかに移すものがあることが、印を忘れた同型の必要十分条件である。
移動した候補とも同じ比較を行えば最小周期を判定できる。
